\section{Introduction}

The paradigm of Large Language Model (LLM) training has undergone a pivotal shift from Supervised Fine-Tuning (SFT)\cite{ouyang2022traininglanguagemodelsfollow} to Reinforcement Learning-based (RL)\cite{Fran_ois_Lavet_2018} alignment. While models such as GPT-4 \cite{Achiam2023GPT4TR} and Claude \cite{TheC3} have demonstrated outstanding reasoning capabilities, their reliance on complex reward models imposes significant computational and scalability bottlenecks. Recently, the release of DeepSeek-R1 popularized Group Relative Policy Optimization (GRPO)\cite{Guo_2025} \cite{shao2024deepseekmathpushinglimitsmathematical}, marking a milestone in training specialized reasoning models without the heavy overhead of a traditional reward model. By leveraging group-relative rewards, GRPO allows the policy model to autonomously distinguish between good and bad reasoning paths within the same prompt context, effectively optimizing for complex system thinking capabilities.

However, as models scale and tasks become more complex, such as Olympiad-level mathematics and complex reasoning chains, vanilla GRPO faces diminishing returns. We identify three fundamental limitations in existing RL pipelines for reasoning:

\textbf{1. Zero Advantage:} In many reasoning tasks, especially at the start or late stages of training, a model may generate rollouts that all share the same outcome, either all correct or all incorrect. In such cases, the mean reward equals the individual reward, and the standard deviation approaches zero, causing the advantage signal to become zero \cite{zhang2026trainlesslearnmore}. This stalls the gradient and wastes expensive compute.

\textbf{2. The Compulsive Correction Fallacy:} Multi-layer \cite{ding2025multilayergrpoenhancingreasoning} or iterative reflection models \cite{madaan2023selfrefineiterativerefinementselffeedback} often suffer from overcorrection. When prompted to reflect, models exhibit a systematic bias towards altering their initial answers, even when correct. This bias degrades the stability of the model outputs and leads to a high rate of false-negative revisions.

\textbf{3. The Noise-Signal Dichotomy:} Traditional RL update paradigms implicitly operate under the assumption that all optimization errors possess equivalent informational value. However, a stochastic error originating from a superficial computational oversight yields significantly less pedagogical utility than a structural error rooted in a fundamental logical misconception. Standard GRPO frameworks currently lack the granular mechanism required to disambiguate these distinct failure modes.

To address these limitations, we propose Semantic-Balanced GRPO (SB-GRPO). Our framework tackles the zero-advantage problem by introducing a Heuristic Confidence Gate, which applies a differentiable switching mechanism to activate the Advantage Calibration mechanism introduced in \cite{nan2025ngrponegativeenhancedgrouprelative}, which we formulate as a Calibrated Advantage ($A_{calib}$), using virtual maximum rewards for uniform batches. To mitigate correction bias and differentiate noise from systematic errors, we integrate a Semantic Error Profiling (SEP) module within a two-layer generation-reflection pipeline. By clustering reasoning paths and constructing Static Balanced Batches, our approach ensures that the model optimizes based on the most highly informative errors rather than trivial computational noise.